\begin{example}[Branches, minimal primes, and normalization scale]
\label{ex:p121-node-cusp-normalized-links}
For the reducible node
\[
 R=\mathbf C[[x,y]]/(xy),
 \qquad \mathfrak m=(x,y),
\]
the normalization is \(\mathbf C[[x]]\times\mathbf C[[y]]\).  Hence
\(\mathcal V_{\mathfrak m}(R)\) consists of two points: the orders on the
two axes, with respective supports \((y)\) and \((x)\).

Minimal primes alone do not record all branches.  The irreducible nodal
local ring
\[
 R'=\left(\mathbf C[x,y]/(y^2-x^2(1+x))\right)_{(x,y)}
\]
is a domain, but its normalization has two points above the closed point.
Indeed, in its fraction field put \(t=y/x\); then
\(x=t^2-1\) and \(y=t(t^2-1)\), and the two centres are \(t=1\) and
\(t=-1\).  Thus \(\mathcal V_{\mathfrak m}(R')\) again has two points,
although \(R'\) has only one minimal prime.

For the cusp
\[
 R''=\mathbf C[[t^2,t^3]]\subset S=\mathbf C[[t]],
 \qquad \mathfrak m''=(t^2,t^3),
\]
there is one branch and
\[
                 \mathcal V_{\mathfrak m''}(R'')
                    =\left\{\tfrac12\operatorname{ord}_t\right\}.
\]
The maximal-ideal normalization on \(S\) instead gives
\(\mathcal V_{(t)}(S)=\{\operatorname{ord}_t\}\).  With the compatible
ideal \(\mathfrak m''S=(t^2)\), however, the upstairs normalized order is
\(\tfrac12\operatorname{ord}_t\), and
\Cref{prop:p126-finite-reduced-curve-link} restricts it exactly to the
point of \(\mathcal V_{\mathfrak m''}(R'')\).  This is the simplest reason
that the change-of-ideal factor in
\eqref{eq:p121-change-normalizing-ideal} cannot be omitted.
\end{example}
